\chapter{텐서와 상대론적 로런츠 힘}

오늘은 스칼라장에서 전기역학으로 넘어갑니다. 앞서 살펴본 내용에서 유용한 유형을 확인했습니다. 하나의 작용으로 장이 입자에 미치는 작용과 입자가 장에 되미치는 작용을 모두 기술할 수 있습니다. 궁극적으로 이 유형을 전자기장에 적용하겠지만, 먼저 계산의 대칭성을 드러낼 만큼 강력한 표기법이 필요합니다.

이 예비 작업은 단순한 장식이 아닙니다. 마이너스 부호, 0, 등호, 벡터 표기에는 그렇지 않으면 말로 되풀이해야 할 개념이 압축되어 있습니다. 텐서 표기법도 상대론에서 같은 역할을 합니다. 위 첨자와 아래 첨자, 축약, 변환 법칙을 갖추고 나면 전기력과 자기력이 하나의 4차원 방정식으로 자연스럽게 모입니다.

\section{4차원 벡터와 첨자 표기의 간결함}

\subsection{위 성분과 아래 성분}

그리스 첨자는 시공간의 네 좌표를 모두 취하고, 라틴 첨자는 세 공간 좌표만 취합니다:
\begin{equation}
\mu,\nu,\rho,\sigma=0,1,2,3,
\qquad
m,n=1,2,3.
\end{equation}
일반적인 반변 4차원 벡터는 \(A^\mu=(A^0,A^1,A^2,A^3)\)입니다.
여기서 사용하는 규약에 따르면,
\begin{equation}
\eta_{\mu\nu}
=
\begin{pmatrix}
-1&0&0&0\\
0&1&0&0\\
0&0&1&0\\
0&0&0&1
\end{pmatrix},
\qquad
A_\mu=\eta_{\mu\nu}A^\nu .
\label{eq:sr6-metric-lowering}
\end{equation}
따라서 첨자를 내리면 시간 성분의 부호가 바뀌고 공간 성분은 그대로입니다:
\begin{equation}
A_0=-A^0,
\qquad
A_m=A^m.
\label{eq:sr6-component-lowering}
\end{equation}
평탄한 시공간에서는 이 연산이 의미하는 바가 전부입니다. 다만 일반상대론에서는 같은 표기법이 더 풍부한 기하학적 의미를 갖게 됩니다.

\begin{figure}[htbp]\centering
\includegraphics[width=0.90\textwidth]{lecture_06_figure_01.png}\caption{민코프스키 계량과 4차원 벡터 첨자를 내리는 성분 규칙.}\lectureframe{00:07:45}\end{figure}

아인슈타인 합산 규약에서는 같은 첨자가 한 번은 위에, 한 번은 아래에 반복될 때 합산합니다. 따라서 축약
\begin{equation}
A_\mu B^\mu
=
-A^0B^0+A^1B^1+A^2B^2+A^3B^3
\label{eq:sr6-vector-contraction}
\end{equation}
은 스칼라입니다. 합이 위 첨자와 아래 첨자가 반복되는 형태가 아니면 합산 기호를 명시하겠습니다--.

\subsection{미분 연산자는 아래 첨자를 갖는다}

네 미분 연산자를 다음과 같이 모아 씁니다.
\begin{equation}
\partial_\mu
\equiv
\frac{\partial}{\partial x^\mu}.
\end{equation}
스칼라장 \(\phi(x)\)에 작용시키면 공변 벡터 \(\partial_\mu\phi\)를 이룹니다. 첨자가 어디에 놓이는지 가장 빠르게 확인하려면 무한소 변위를 살펴보면 됩니다:
\begin{equation}
d\phi
=
\frac{\partial\phi}{\partial x^\mu}\,dx^\mu
=
\partial_\mu\phi\,dx^\mu.
\label{eq:sr6-total-differential}
\end{equation}
좌변은 이웃한 사건에서의 스칼라값 차이이므로 스칼라입니다. \(dx^\mu\)가 반변이므로, 여기에 축약되는 대상은 공변이어야 합니다.

\begin{classroomqa}\audiencequestion{\(x^\mu\)에 위 첨자가 있는데, \(x^\mu\)에 대해 미분하면 왜 아래 첨자가 생기나요?}\lecturerresponse{미소 변화량 \(d\phi=\partial_\mu\phi\,dx^\mu\)은 스칼라이기 때문입니다. 반변 변위와 축약해 스칼라를 만들려면 공변 기울기가 필요합니다. 따라서 미분 첨자의 위치는 분수의 글자 모양이 아니라 변환 법칙으로 정해집니다. \lecturetimestamp{00:10:10}}\end{classroomqa}

\begin{figure}[htbp]\centering
\includegraphics[width=0.90\textwidth]{lecture_06_figure_02.png}\caption{\(d\phi\)이 스칼라 미분량이므로 \(\partial_\mu\phi\)은 공변 벡터임을 알 수 있다.}\lectureframe{00:11:05}\end{figure}

같은 규칙으로 또 다른 스칼라를 얻습니다. \(B^\mu(x)\)가 4차원 벡터장이라면,
\begin{equation}
\partial_\mu B^\mu
\label{eq:sr6-four-divergence}
\end{equation}
은 그 4차원 발산입니다. 시간 성분에 대한 시간 미분과 대응하는 공간 성분들에 대한 세 공간 미분을 하나로 묶습니다.

\begin{figure}[htbp]\centering
\includegraphics[width=0.90\textwidth]{lecture_06_figure_03.png}\caption{두 스칼라 축약: \(A_\mu B^\mu\)와 4차원 발산 \(\partial_\mu B^\mu(x)\).}\lectureframe{00:14:00}\end{figure}

이 모든 규칙에는 3차원 대응물이 있습니다. 유클리드 3차원 공간에서 계량은 항등행렬이므로 라틴 첨자를 올리거나 내려도 달라지는 것이 없습니다. 그러므로 편리한 경우에는 공간 첨자를 위아래로 바꿔 쓰되, 시공간 첨자는 계속 엄밀히 구별하겠습니다.

\section{행렬로 나타낸 로런츠 변환}

스칼라와 벡터는 변환하는 방식으로 정의됩니다. 스칼라는 값이 그대로 유지되지만, 반변 4차원 벡터는 다음을 따릅니다.
\begin{equation}
(A^\mu)'=L^\mu{}_\nu A^\nu.
\label{eq:sr6-vector-transform}
\end{equation}
눈에 보이는 첨자 \(\mu\)는 양변에 나타나는 자유 첨자입니다. 반복 첨자 \(\nu\)는 합산되며 최종 성분 방정식에는 남지 않습니다.

여기서 \emph{로런츠 변환}에는 공간 회전과 부스트가 모두 포함됩니다. 부스트는 어느 축을 따라서든 일어날 수 있으며, 부스트와 회전을 합성하면 일반적인 변환을 만들 수 있습니다. 공간 회전 자체는 시간 좌표를 바꾸지 않습니다.

\begin{classroomqa}\audiencequestion{공간을 회전하면 시간은 전혀 영향을 받지 않나요?}\lecturerresponse{그렇습니다. 일반적인 공간 회전에서는 시간 성분이 변하지 않습니다. 시간과 공간을 섞는 변환들과 함께 이런 회전도 로런츠 군에 포함됩니다. \lecturetimestamp{00:16:27}}\end{classroomqa}

\(v\)의 속력으로 \(x\)축 방향으로 부스트할 때, \(c=1\)이면
\begin{equation}
L^\mu{}_\nu
=
\begin{pmatrix}
\gamma&-\gamma v&0&0\\
-\gamma v&\gamma&0&0\\
0&0&1&0\\
0&0&0&1
\end{pmatrix},
\qquad
\gamma=\frac{1}{\sqrt{1-v^2}}.
\label{eq:sr6-boost-matrix}
\end{equation}
\((t,x,y,z)^{\mathsf T}\)를 곱하면
\begin{equation}
\begin{aligned}
t'&=\gamma(t-vx),\\
x'&=\gamma(x-vt),\\
y'&=y,\\
z'&=z.
\end{aligned}
\end{equation}
행렬을 칠판에 완성할 때, 수업에서는 영인 항들과 단위값을 갖는 두 항을 채워 넣습니다.

\begin{figure}[htbp]\centering
\includegraphics[width=0.90\textwidth]{lecture_06_figure_04.png}\caption{열벡터 \((t,x,y,z)^{\mathsf T}\)에 작용하는 \(x\) 방향 부스트 행렬.}\lectureframe{00:21:35}\end{figure}

반면 \(y\)--\(z\) 평면에서 각도 \(\theta\)만큼 회전하면 다음과 같습니다.
\begin{equation}
R_{yz}(\theta)
=
\begin{pmatrix}
1&0&0&0\\
0&1&0&0\\
0&0&\cos\theta&\sin\theta\\
0&0&-\sin\theta&\cos\theta
\end{pmatrix}.
\label{eq:sr6-yz-rotation}
\end{equation}
이는 다음을 뜻합니다.
\begin{equation}
y'=y\cos\theta+z\sin\theta,
\qquad
z'=-y\sin\theta+z\cos\theta,
\end{equation}
한편 \(t\)와 \(x\)는 그대로 유지됩니다. \eqref{eq:sr6-boost-matrix}와 \eqref{eq:sr6-yz-rotation} 같은 행렬의 곱으로 더 일반적인 로런츠 변환을 만들 수 있습니다.

\begin{figure}[htbp]\centering
\includegraphics[width=0.90\textwidth]{lecture_06_figure_05.png}\caption{\(y\)--\(z\) 평면에서 일어나는 순수 회전. \(t\)와 \(x\)는 변하지 않는다.}\lectureframe{00:23:10}\end{figure}

공변 벡터는 첨자의 위치가 반대인 행렬을 사용해 변환합니다:
\begin{equation}
(A_\mu)'=M_\mu{}^\nu A_\nu.
\label{eq:sr6-covector-transform}
\end{equation}
칠판 앞에서 청중이 지적한 대로, 잠시 반복된 자유 첨자가 있습니다. 좌변의 합산되지 않는 \(\mu\)는 우변에도 합산되지 않는 첨자로 나타나야 합니다. 두 행렬은 같은 물리적 기준계 변환을 나타내며 다음 관계를 만족합니다.
\begin{equation}
M=\eta L\eta,
\qquad
\eta^{-1}=\eta.
\label{eq:sr6-M-L-relation}
\end{equation}
마지막 항등식은 \(\eta\)의 모든 대각 성분이 자기 자신의 역원이라는 사실에서 바로 따릅니다.

\begin{figure}[htbp]\centering
\includegraphics[width=0.90\textwidth]{lecture_06_figure_06.png}\caption{반변 변환 법칙과 공변 변환 법칙, 그리고 둘 사이의 관계 \(M=\eta L\eta\).}\lectureframe{00:25:20}\end{figure}

이제 핵심 개념을 사용할 수 있습니다. 4차원 벡터란 그 성분이 시공간 좌표와 나란히 변환하는 대상입니다.

\section{텐서: 첨자마다 하나씩 적용되는 변환}

\subsection{벡터의 곱에서 정의까지}

스칼라는 계수 0 텐서이고, 벡터는 계수 1 텐서이며, 첨자가 둘인 대상은 계수 2 텐서입니다. 두 벡터에서 시작해 다음을 만듭니다.
\begin{equation}
T^{\mu\nu}=A^\mu B^\nu.
\label{eq:sr6-simple-tensor}
\end{equation}
성분은 \(4\times4=16\)개입니다. 성분에는 \(A^0B^0\)와 \(A^0B^1\)부터 \(A^3B^3\)까지가 있습니다.
기호 \(T^{\mu\nu}\)는 성분 하나가 아니라 이 배열 전체를 나타냅니다.

각 벡터를 따로 변환하면
\begin{align}
(A^\mu)'&=L^\mu{}_\sigma A^\sigma,\\
(B^\nu)'&=L^\nu{}_\tau B^\tau,
\end{align}
을 얻고, 따라서
\begin{equation}
(T^{\mu\nu})'
=
L^\mu{}_\sigma L^\nu{}_\tau T^{\sigma\tau}.
\label{eq:sr6-rank-two-transform}
\end{equation}
이 성립합니다.
텐서의 각 첨자에는 각자의 로런츠 행렬이 적용됩니다. 더미 첨자 \(\sigma\)와 \(\tau\)는 각 쌍 안에서 일관성을 유지하는 한 임의로 이름을 바꿀 수 있습니다.

\begin{classroomqa}\audiencequestion{이 재배열은 행렬과 벡터 성분이 서로 교환된다고 가정하나요?}\lecturerresponse{여기서 다루는 것은 서로 교환하지 않는 양자 연산자가 아니라 보통의 고전적인 수이므로 서로 교환됩니다. 게다가 같은 로런츠 행렬이 두 첨자 모두에 작용합니다. 이 계산에는 순서 문제가 없습니다. \lecturetimestamp{00:31:56}}\end{classroomqa}

이 규칙은 계수 3으로 바로 확장됩니다:
\begin{equation}
(T^{\mu\nu\lambda})'
=
L^\mu{}_\sigma
L^\nu{}_\tau
L^\lambda{}_\kappa
T^{\sigma\tau\kappa},
\label{eq:sr6-rank-three-transform}
\end{equation}
다른 계수에서도 마찬가지입니다. 특정하게 벡터들을 곱해 표현할 수 있느냐가 아니라 이 변환 법칙이 텐서를 정의합니다.
예를 들어,
\begin{equation}
A^\mu B^\nu+C^\mu D^\nu
\end{equation}
도 텐서입니다. 일반적으로 두 벡터의 곱 하나로 나타낼 수 없어도 그렇습니다.

\begin{figure}[htbp]\centering
\includegraphics[width=0.90\textwidth]{lecture_06_figure_07.png}\caption{벡터의 곱에서 계수 2 및 계수 3 텐서의 변환 법칙으로 이어지는 칠판 전개.}\lectureframe{00:33:20}\end{figure}

\begin{classroomqa}\audiencequestion{계수 2 텐서의 행렬 곱을 성분별로 쓰면 어떤 뜻인가요?}\lecturerresponse{출력 첨자 둘을 고정하고 숨은 첨자 둘에 대해 합산합니다. 예를 들면 다음과 같습니다.
\[
(T^{31})'=L^3{}_\sigma L^1{}_\tau T^{\sigma\tau}.
\]
\(\sigma\)와 \(\tau\)가 각각 네 값을 취하므로 이 출력 성분 하나에는 항이 16개 들어 있습니다. \lecturetimestamp{00:34:35}}\end{classroomqa}

\begin{figure}[htbp]\centering
\includegraphics[width=0.90\textwidth]{lecture_06_figure_08.png}\caption{\((T^{31})'\)를 전개하면 서로 독립인 두 합산이 드러난다.}\lectureframe{00:35:55}\end{figure}

텐서 방정식에는 결정적인 장점이 있습니다. 두 텐서의 성분이 한 기준계에서 같으면, 변환 법칙에 따라 모든 기준계에서 두 텐서가 같습니다. 달리 말해, 한 기준계에서 텐서의 모든 성분이 영이면 모든 기준계에서 텐서 전체가 영입니다. 성분 하나만 영이라고 해서 충분한 것은 아닙니다. 모든 성분이 영이어야 합니다. 따라서 텐서 사이에 쓴 방정식은 한 관성계에서 다른 관성계로 옮겨도 참입니다.

\subsection{텐서 첨자의 올림과 내림}

텐서의 첨자를 내리는 방법은 벡터의 첨자를 내리는 방법과 정확히 같습니다:
\begin{equation}
T^\mu{}_\nu=T^{\mu\sigma}\eta_{\sigma\nu}.
\label{eq:sr6-tensor-lowering}
\end{equation}
\eqref{eq:sr6-metric-lowering}의 대각 계량에서는 시간 첨자를 내릴 때마다 마이너스 부호가 하나씩 붙습니다. 따라서
\begin{equation}
T_{00}=T^{00},
\qquad
T_{01}=-T^{01}.
\label{eq:sr6-tensor-component-signs}
\end{equation}
첫 번째 등식에는 부호 변화가 두 번 있고, 두 번째에는 한 번만 있습니다.

\begin{figure}[htbp]\centering
\includegraphics[width=0.90\textwidth]{lecture_06_figure_09.png}\caption{텐서 첨자 두 개를 내리는 예: 시간 첨자가 둘이면 마이너스 부호가 두 개, 시간 첨자가 하나이면 하나가 붙는다.}\lectureframe{00:41:00}\end{figure}

\begin{classroomqa}\audiencequestion{공변 형태와 반변 형태를 모두 유지하면 어떤 점이 유용한가요? 표기를 간결하게 하는 것 말고도 이유가 있나요?}\lecturerresponse{현재 계산에서 얻는 직접적인 이점은 축약이 간결해지고 민코프스키 부호를 자동으로 관리할 수 있다는 점입니다. 벡터와 쌍대벡터 사이에는 더 깊은 기하학적 구별이 있지만, 이는 일반상대론에서 다루는 편이 낫습니다. 평탄한 시공간에서는 첨자 규칙이 매우 효율적이어서 보통 이를 직접 사용합니다. \lecturetimestamp{00:41:26}}\end{classroomqa}

\begin{classroomqa}\audiencequestion{텐서에는 항상 성분이 네 개 있나요, 아니면 상대론에서만 그런가요?}\lecturerresponse{각 첨자가 취하는 값의 범위는 공간의 차원으로 정해집니다. 4차원 시공간에서 계수 1 텐서의 성분은 네 개이고 계수 2 텐서의 성분은 16개입니다. 3차원에서는 각각 세 개와 아홉 개입니다. 텐서의 변환 규칙 자체는 차원에 의존하지 않습니다. \lecturetimestamp{00:42:55}}\end{classroomqa}

\begin{classroomqa}\audiencequestion{공변 대상과 반변 대상을 더 기하학적으로 설명할 수 있나요?}\lecturerresponse{그렇습니다. 기하학적 해석이 있지만 지금 이를 따라가면 전기역학을 더 미루게 됩니다. 이 강의에서는 변환 및 첨자 조작 규칙으로 충분합니다. 기하학적 구별은 일반상대론에서 피할 수 없는 주제가 됩니다. \lecturetimestamp{00:43:51}}\end{classroomqa}

\section{대칭성, 반대칭성, 전자기 텐서}

일반적인 계수 2 텐서에서는 첨자를 맞바꾸면 성분이 달라집니다:
\begin{equation}
T^{\mu\nu}\neq T^{\nu\mu}.
\end{equation}
실제로 \(A^\mu B^\nu\)와 \(A^\nu B^\mu\)는 일반적으로 서로 다르므로, 특정 성분 \(A^0B^1\)와 \(A^1B^0\)도 같을 필요가 없습니다.

중요한 특수한 경우가 두 가지 있습니다. 대칭 텐서는 다음을 만족합니다.
\begin{equation}
T^{\mu\nu}=T^{\nu\mu},
\end{equation}
구체적인 대칭 조합은 다음과 같습니다.
\begin{equation}
A^\mu B^\nu+A^\nu B^\mu.
\end{equation}
반대칭 텐서는 다음을 만족합니다.
\begin{equation}
F^{\mu\nu}=-F^{\nu\mu},
\label{eq:sr6-antisymmetry}
\end{equation}
다음 조합도 마찬가지입니다.
\begin{equation}
A^\mu B^\nu-A^\nu B^\mu.
\end{equation}
\eqref{eq:sr6-antisymmetry}에서 두 첨자를 같게 놓으면
\begin{equation}
F^{00}=-F^{00},
\qquad
F^{00}=0,
\end{equation}
을 얻으며, 같은 논리로 모든 대각 성분이 사라집니다. 따라서 \(4\times4\) 반대칭 텐서에는 서로 독립인 성분이 여섯 개 있습니다.

\begin{figure}[htbp]
\centering
\includegraphics[width=0.90\textwidth]{lecture_06_figure_10.png}
\caption{칠판에 비교해 놓은 대칭 이계 텐서와 반대칭 이계 텐서.}
\lectureframe{00:47:20}
\end{figure}

이 여섯 성분은 전기장의 세 성분과 자기장의 세 성분에 정확히 대응합니다. 칠판에 적은 부호 규약을 쓰면,
\begin{equation}
F^{\mu\nu}
=
\begin{pmatrix}
0&-E_1&-E_2&-E_3\\
E_1&0&-B_3&B_2\\
E_2&B_3&0&-B_1\\
E_3&-B_2&B_1&0
\end{pmatrix}.
\label{eq:sr6-em-field-tensor}
\end{equation}
대각선을 기준으로 반사하면 모든 부호가 바뀝니다. 따라서 전기장과 자기장은 상대론적으로 서로 독립된 대상이 아닙니다. 부스트는 시간과 공간을 섞듯이 두 장을 섞을 수 있습니다. 한 기준계에서 순수하게 전기장만 있던 장이 다른 기준계에서는 자기장 성분을 가질 수 있습니다.

\begin{figure}[htbp]\centering
\includegraphics[width=0.90\textwidth]{lecture_06_figure_11.png}\caption{서로 독립인 여섯 성분 \(\mathbf E\)와 \(\mathbf B\)로 이루어진 완성된 반대칭 배열.}\lectureframe{00:50:20}\end{figure}

\eqref{eq:sr6-em-field-tensor}를 적은 칠판을 그대로 둔 채 논의는 입자 동역학으로 넘어갑니다. 이제 기호 \(A_\mu\)를 다시 사용하지만, 앞선 예에서 나온 일반적인 벡터가 아니라 전자기 4차원 퍼텐셜을 뜻합니다.

\section{하전 입자의 작용}

출발점은 두 가지입니다. 장 자체의 동역학을 살펴보면 맥스웰 방정식으로 이어지고, 주어진 장 속에서 움직이는 하전 입자를 살펴보면 로런츠 힘으로 이어집니다. 여기서는 입자부터 시작하겠습니다.

저속도에서 목표로 삼는 식은
\begin{equation}
m\mathbf a
=
e\left(\mathbf E+\mathbf v\times\mathbf B\right),
\label{eq:sr6-nonrel-lorentz-force}
\end{equation}
입니다. 가우스 단위계에서는 자기력 항에 계수 \(1/c\)가 곱해집니다. \(c=1\)로 놓겠습니다. \eqref{eq:sr6-nonrel-lorentz-force}의 두 항은 3차원 벡터 표기로는 서로 다르게 보이지만, 상대론에서는 하나의 대상의 일부임이 드러나야 합니다.

\begin{figure}[htbp]\centering
\includegraphics[width=0.90\textwidth]{lecture_06_figure_12.png}\caption{불변 작용이 재현해야 하는 결과인 비상대론적 로런츠 힘.}\lectureframe{00:54:00}\end{figure}

\subsection{자유 운동과 벡터 퍼텐셜}

자유 상대론적 입자에 대해서는,
\begin{equation}
S_{\mathrm{free}}
=
-m\int_1^2 d\tau
=
\int_{t_1}^{t_2}dt\,
\left[-m\sqrt{1-\dot{\mathbf x}^{\,2}}\right],
\label{eq:sr6-free-action}
\end{equation}
이며,
\begin{equation}
\dot{\mathbf x}^{\,2}
=
(\dot x^1)^2+(\dot x^2)^2+(\dot x^3)^2.
\end{equation}
입니다. 대괄호 안의 양은 자유 입자의 라그랑지안입니다. 운동량과 에너지는 앞서 구했습니다. 여기서는 작용을 유지하고 장에 의존하는 항을 무엇을 더할 수 있는지 살펴보겠습니다.

\begin{figure}[htbp]\centering
\includegraphics[width=0.90\textwidth]{lecture_06_figure_13.png}\caption{세계선에 따라 누적되는 고유시간으로 나타낸 자유 입자의 작용.}\lectureframe{00:55:10}\end{figure}

전자기학의 기본 변수는 4차원 벡터 퍼텐셜 \(A_\mu(x)\)입니다. 여기서 전기장과 자기장을 유도하겠습니다. 세계선의 작은 구간 \(dx^\mu\)에 대해 축약
\begin{equation}
A_\mu(x)\,dx^\mu
\end{equation}
은 스칼라입니다. 한 끝점에서 다른 끝점까지 이러한 스칼라 미소량을 더하면 로런츠 불변인 상호작용을 얻습니다.
\begin{equation}
S_{\mathrm{int}}
=
-e\int_1^2 A_\mu(x)\,dx^\mu.
\label{eq:sr6-interaction-action}
\end{equation}
계수 \(e\)는 전하입니다. 전체 부호는 역사적으로 정해진 전하 규약에 따릅니다. 규약을 바꾸면 \(e\)에 붙는 부호가 바뀌지만 이론의 내용은 달라지지 않습니다.

\begin{figure}[htbp]\centering
\includegraphics[width=0.90\textwidth]{lecture_06_figure_14.png}\caption{세계선 구간과 4차원 퍼텐셜의 축약으로 스칼라 미소량 \(A_\mu dx^\mu\)를 만든다.}\lectureframe{00:59:20}\end{figure}

\subsection{보통 시간으로 쓴 형태}

익숙한 오일러--라그랑주 방정식을 사용하려면 \eqref{eq:sr6-interaction-action}에 \(dt\)를 곱하고 나누어야 합니다. 다음 관계가 성립하므로
\begin{equation}
\frac{dx^0}{dt}=1,
\qquad
\frac{dx^m}{dt}=\dot x^m,
\end{equation}
전체 작용은 다음과 같습니다.
\begin{equation}
S
=
\int dt\,
\left[
-m\sqrt{1-\dot{\mathbf x}^{\,2}}
-eA_0(x,t)
-e\dot x^m A_m(x,t)
\right].
\label{eq:sr6-complete-particle-action}
\end{equation}
따라서
\begin{equation}
\boxed{
L
=
-m\sqrt{1-\dot{\mathbf x}^{\,2}}
-eA_0
-e\dot x^mA_m
}.
\label{eq:sr6-charged-lagrangian}
\end{equation}
입니다. 지금은 \(A_\mu(x,t)\)가 주어진 외부장입니다. 변분하는 것은 장이 아니라 세계선입니다.

\begin{figure}[htbp]\centering
\includegraphics[width=0.90\textwidth]{lecture_06_figure_15.png}\caption{보통 시간의 라그랑지안으로 다시 쓴 세계선의 불변 결합.}\lectureframe{01:05:20}\end{figure}

\section{로런츠 힘을 직접 계산하기}

공간 성분 \(x^m\)에 대한 오일러--라그랑주 방정식은 다음과 같습니다.
\begin{equation}
\frac{d}{dt}\frac{\partial L}{\partial\dot x^m}
=
\frac{\partial L}{\partial x^m}.
\label{eq:sr6-EL}
\end{equation}
속도에 대해 미분하면
\begin{equation}
\frac{\partial L}{\partial\dot x^m}
=
\frac{m\dot x_m}{\sqrt{1-\dot{\mathbf x}^{\,2}}}
-eA_m.
\label{eq:sr6-canonical-momentum}
\end{equation}
을 얻습니다.
불행히도 문자 \(m\)는 질량과 공간 첨자 양쪽에 쓰입니다. 문맥으로 구별할 수 있습니다. \eqref{eq:sr6-canonical-momentum}에서 첫 번째 \(m\)는 질량이고, 아래 첨자 \(m\)는 세 공간 성분 중 하나를 나타냅니다.

\begin{figure}[htbp]\centering
\includegraphics[width=0.90\textwidth]{lecture_06_figure_16.png}\caption{운동량의 미분에는 역학적 항과 벡터 퍼텐셜 항이 모두 들어간다.}\lectureframe{01:09:20}\end{figure}

\begin{classroomqa}\audiencequestion{좌변을 처음 계산할 때 항 하나가 빠진 것 아닌가요?}\lecturerresponse{그렇습니다. 자유 입자 항만 미분하면 \(-e\dot x^mA_m\)의 속도 의존성을 빠뜨립니다. 전체 미분은 \(m\dot x_m/\sqrt{1-\dot{\mathbf x}^{\,2}}-eA_m\)입니다. 이 보정항은 시간 미분을 취할 때 전기력과 자기력 양쪽에 일부를 보태므로 꼭 필요합니다. \lecturetimestamp{01:08:18}}\end{classroomqa}

\eqref{eq:sr6-EL}의 우변은
\begin{equation}
\frac{\partial L}{\partial x^m}
=
-e\,\partial_mA_0
-e\,\dot x^n\partial_mA_n.
\label{eq:sr6-position-derivative}
\end{equation}
입니다. 따라서
\begin{equation}
\frac{d}{dt}
\left(
\frac{m\dot x_m}{\sqrt{1-\dot{\mathbf x}^{\,2}}}
-eA_m
\right)
=
-e\,\partial_mA_0
-e\,\dot x^n\partial_mA_n.
\label{eq:sr6-EL-expanded}
\end{equation}

\begin{classroomqa}\audiencequestion{이 계산에서 공변 공간 성분과 반변 공간 성분을 서로 바꾸어 쓴 것 아닌가요?}\lecturerresponse{첨자를 올리고 내리기는 했지만, 라틴 공간 첨자에 해당하는 계량의 유클리드 공간 블록은 항등행렬입니다. 따라서 \(x_m=x^m\)이고 \(A_m=A^m\)입니다. 부호가 바뀌는 것은 시간 첨자뿐입니다. \lecturetimestamp{01:09:29}}\end{classroomqa}

\begin{classroomqa}\audiencequestion{반복되는 라틴 첨자도 합산하나요?}\lecturerresponse{그렇습니다. 반복되는 라틴 첨자는 \(1\)부터 \(3\)까지 합산하고, 반복되는 그리스 첨자는 \(0\)부터 \(3\)까지 합산합니다. 공간 방정식 하나를 표시하는 \(m\)처럼 하나만 있는 첨자는 자유 첨자이므로 합산하지 않습니다. 질량을 나타내는 기호는 첨자가 아닙니다. \lecturetimestamp{01:10:00}}\end{classroomqa}

\begin{figure}[htbp]\centering
\includegraphics[width=0.90\textwidth]{lecture_06_figure_17.png}\caption{벡터 퍼텐셜의 전미분을 전개하기 전 오일러--라그랑주 방정식.}\lectureframe{01:14:20}\end{figure}

\subsection{움직이는 세계선을 따른 연쇄 법칙}

시간에 따라 변하지 않는 장이라도 입자가 서로 다른 위치를 지나면서 그 경로를 따라 변합니다. 전미분은 다음과 같습니다.
\begin{equation}
\frac{dA_m(x(t),t)}{dt}
=
\partial_0A_m+\dot x^n\partial_nA_m.
\label{eq:sr6-convective-derivative}
\end{equation}
첫째 항은 시간에 대한 명시적 의존성을 나타내고, 둘째 항은 \(x^n(t)\)를 통한 암묵적 시간 의존성을 나타냅니다.

\begin{figure}[htbp]\centering
\includegraphics[width=0.90\textwidth]{lecture_06_figure_18.png}\caption{벡터 퍼텐셜은 시간에 명시적으로 의존하고, 입자가 공간을 지나 움직이는 데 따라 암묵적으로도 변한다.}\lectureframe{01:17:20}\end{figure}

\eqref{eq:sr6-convective-derivative}를 \eqref{eq:sr6-EL-expanded}에 대입하고 장 항을 우변으로 옮기면
\begin{equation}
\begin{aligned}
m\frac{d}{dt}
\left(
\frac{\dot x_m}{\sqrt{1-\dot{\mathbf x}^{\,2}}}
\right)
&=
e\left(\partial_0A_m-\partial_mA_0\right)
\\
&\quad +e\dot x^n
\left(\partial_nA_m-\partial_mA_n\right).
\end{aligned}
\label{eq:sr6-grouped-force}
\end{equation}
마침내 계산을 정리하는 원리를 찾았습니다. 첫 번째 괄호 안의 항은 속도와 무관하고, 두 번째 항은 속도에 선형입니다. 전기력과 자기력의 구분이 바로 이와 같습니다.

\begin{figure}[htbp]\centering
\includegraphics[width=0.90\textwidth]{lecture_06_figure_19.png}\caption{속도와 무관한 전기력 항과 속도에 의존하는 반대칭 항으로 묶은 힘 항들.}\lectureframe{01:23:40}\end{figure}

유도 과정의 규약에 따라 전기장 성분을 다음과 같이 정의합니다.
\begin{equation}
E_m=\partial_0A_m-\partial_mA_0.
\label{eq:sr6-electric-from-potential}
\end{equation}
공간 반대칭 조합
\begin{equation}
C_{mn}
\equiv
\partial_nA_m-\partial_mA_n
\label{eq:sr6-spatial-curl-tensor}
\end{equation}
에는 벡터 퍼텐셜의 회전의 서로 독립인 세 성분이 들어 있습니다. 방향은 다음을 만족하도록 정합니다.
\begin{equation}
\dot x^nC_{mn}
=
\left(\mathbf v\times\mathbf B\right)_m.
\end{equation}
따라서 방정식 \eqref{eq:sr6-grouped-force}은 \eqref{eq:sr6-nonrel-lorentz-force}의 상대론적 운동량 형태입니다. 시간 미분 아래에 있는 인자는 단순한 뉴턴 역학의 속도가 아니라 단위 질량당 상대론적 운동량입니다.

직접 손으로 유도 과정을 다시 해볼 만합니다. \eqref{eq:sr6-charged-lagrangian}에서 시작해 \eqref{eq:sr6-EL}의 양변을 계산하고, 연쇄 법칙을 적용한 뒤, 속도 의존성에 따라 항을 묶어 보세요. 이 연습으로 변분 역학, 다변수 미적분학, 전기역학을 함께 익힐 수 있습니다.

\begin{classroomqa}\audiencequestion{전자기 상호작용으로 \(-e\int A_\mu dx^\mu\)를 택하는 이유는 무엇인가요?}\lecturerresponse{궁극적으로 결합의 형태는 실험으로 결정되지만, 두 가지 구조적 조건을 보면 이것이 자연스러운 첫 번째 후보입니다. 작용은 궤적을 따라 국소적인 미소량을 모아 만들며, 각 미소량은 스칼라여야 합니다. 스칼라 \(dx^\mu dx_\mu=d\tau^2\)는 이미 자유 입자의 작용으로 표현되어 있습니다. \(dx^\mu\)를 장 \(A_\mu\)와 축약하는 것이 장에 의존하는 가장 단순한 새 스칼라입니다. 더 복잡한 스칼라도 쓸 수 있으며, 뒤에서 게이지 불변성이 선택의 폭을 더 좁혀 줄 것입니다. \lecturetimestamp{01:25:34}}\end{classroomqa}

\begin{classroomqa}\audiencequestion{최소 작용 원리로 맥스웰 방정식도 유도할 수 있나요?}\lecturerresponse{그렇습니다. 지금 다루는 작용은 주어진 전자기장 속 입자의 운동을 기술합니다. 맥스웰 방정식을 얻으려면 장 자체의 작용을 더하고, 앞서 스칼라장에서 했듯이 장 자체를 변분해야 합니다. \lecturetimestamp{01:28:58}}\end{classroomqa}

\begin{classroomqa}\audiencequestion{묶어 쓴 힘 방정식에서 합산은 어디에 있나요?}\lecturerresponse{자유 첨자 \(m\)는 세 성분 방정식 중 하나를 표시하며 합산하지 않습니다. \(\dot x^n(\partial_nA_m-\partial_mA_n)\)에서 반복되는 \(n\)는 \(1,2,3\)에 대해 합산합니다. 아인슈타인 표기법 없이 쓰면 이 항에 명시적인 \(\sum_{n=1}^{3}\)가 붙습니다. \lecturetimestamp{01:29:26}}\end{classroomqa}

\begin{classroomqa}\audiencequestion{통상적인 공식을 미리 가정하는 대신 \(\partial_0A_m-\partial_mA_0\)를 전기장으로 식별하는 이유는 무엇인가요?}\lecturerresponse{작용을 통해 입자의 속도에는 의존하지 않고 위치와 시간에 의존하는 힘 항을 얻었습니다. 이것이 바로 전기장의 작동상 역할입니다. 나머지 반대칭 항은 속도에 곱해지므로 자기력의 역할을 합니다. 이 계산은 별도의 정의를 미리 대입하는 대신 힘 법칙에서 장을 발견하는 과정으로 볼 수 있습니다. \lecturetimestamp{01:31:05}}\end{classroomqa}

\section{명시적 공변성을 갖는 힘 법칙}

\eqref{eq:sr6-grouped-force}의 좌변은 아직 4차원 형태가 분명히 드러나도록 쓰이지 않았습니다. 4차원 속도를 도입하면
\begin{equation}
u^\mu=\frac{dx^\mu}{d\tau}.
\label{eq:sr6-four-velocity}
\end{equation}
공간 성분은 다음을 만족합니다.
\begin{equation}
u^m
=
\frac{dx^m}{d\tau}
=
\frac{\dot x^m}{\sqrt{1-\dot{\mathbf x}^{\,2}}}.
\end{equation}
공간 방정식에 \(dt/d\tau\)를 곱하면 \(d/dt\)는 \(d/d\tau\)로 바뀝니다:
\begin{equation}
\frac{dt}{d\tau}\frac{d}{dt}
=
\frac{d}{d\tau}.
\end{equation}
따라서 좌변은 고유 가속도의 공간 성분이 됩니다.
\begin{equation}
m\frac{d^2x^m}{d\tau^2}.
\end{equation}

우변도 시간과 공간을 나누지 않고 쓸 수 있습니다.
\begin{equation}
m\frac{d^2x^m}{d\tau^2}
=
e\left(\partial_\nu A_m-\partial_mA_\nu\right)
\frac{dx^\nu}{d\tau}.
\label{eq:sr6-spatial-covariant-force}
\end{equation}
\(\nu=0\)에서는 전기장 괄호를 재현하고, \(\nu=1,2,3\)에서는 자기장 괄호를 재현합니다.

이제 양변 모두 4차원 벡터의 변환 성질을 갖습니다. 따라서 세 공간 방정식은 다음 방정식의 성분입니다.
\begin{equation}
\boxed{
m\frac{d^2x^\mu}{d\tau^2}
=
e\left(\partial_\nu A^\mu-\partial^\mu A_\nu\right)u^\nu
}.
\label{eq:sr6-covariant-force}
\end{equation}
다음과 같이 정의하면
\begin{equation}
F_{\mu\nu}
\equiv
\partial_\mu A_\nu-\partial_\nu A_\mu,
\label{eq:sr6-F-from-A}
\end{equation}
이 강의에서 사용한 첨자 순서와 상호작용 부호에 따라 동등한 간결한 형태는 다음과 같습니다.
\begin{equation}
m\frac{du^\mu}{d\tau}
=
-eF^\mu{}_\nu u^\nu.
\label{eq:sr6-covariant-force-F}
\end{equation}
\editorialnote{전기역학 교재마다 계량의 부호 규약, \(A_\mu\)의 정의, 전하, \(F_{\mu\nu}\) 사이에 부호를 다르게 배분합니다. \eqref{eq:sr6-covariant-force}과 \eqref{eq:sr6-covariant-force-F}은 작용 \(-e\int A_\mu dx^\mu\)에서 얻은 부호와 첨자 순서를 유지합니다. 불변인 내용은 퍼텐셜의 반대칭화된 미분입니다.}

\begin{figure}[htbp]\centering
\includegraphics[width=0.90\textwidth]{lecture_06_figure_20.png}\caption{고유시간을 사용해 세 공간 방정식을 하나의 4차원 벡터 방정식으로 모은다.}\lectureframe{01:42:40}\end{figure}

\begin{classroomqa}\audiencequestion{네 번째인 시간 성분 방정식을 더해도 되는 이유는 무엇인가요? \(d^2x^0/d\tau^2\)가 단순히 0인가요?}\lecturerresponse{일반적으로 0은 아닙니다. \(x^0=t\)이지만 입자의 속도가 변하면 \(dt/d\tau=\gamma\)도 변합니다. 결정적인 논거는 공변성입니다. 스칼라 작용은 로런츠 불변 방정식을 보장합니다. 공간 성분이 하나의 4차원 벡터 방정식의 성분임을 확인하면, 로런츠 변환은 이들을 시간 성분과 섞습니다. 그러므로 나머지 세 방정식에 네 번째 방정식도 함께 있어야 합니다. \lecturetimestamp{01:42:40}}\end{classroomqa}

\subsection{네 번째 성분의 의미}

세 공간 성분은 상대론적으로 다음을 나타냅니다.
\begin{equation}
\frac{d\mathbf p}{dt}=\mathbf F.
\end{equation}
에너지와 운동량은 그 자체로 4차원 벡터의 성분이므로, 시간 성분은 이에 대응하는 에너지 균형을 나타내야 합니다. 3차원 벡터로 쓰면,
\begin{equation}
\frac{dK}{dt}=\mathbf F\cdot\mathbf v.
\label{eq:sr6-work-energy}
\end{equation}
이는 장이 입자에 한 일의 변화율을 나타냅니다. 자기력은 \(\mathbf v\)에 수직이므로 일을 하지 않습니다. 전기력은 입자의 에너지를 바꿉니다.

\begin{classroomqa}\audiencequestion{공변성으로 네 번째 방정식을 형식적으로 얻는다면, 그 방정식이 물리적 의미를 갖는다는 것을 어떻게 알 수 있나요?}\lecturerresponse{처음 세 성분은 힘이 운동량을 바꾼다고 말합니다. 네 번째 성분은 일이 에너지를 바꾼다고 말합니다: \(dK/dt=\mathbf F\cdot\mathbf v\). 에너지와 운동량이 하나의 4차원 벡터를 이루므로, 장에 대한 반응 역시 네 성분으로 이루어진 하나의 진술이어야 합니다. \lecturetimestamp{01:45:20}}\end{classroomqa}

\begin{figure}[htbp]\centering
\includegraphics[width=0.90\textwidth]{lecture_06_figure_21.png}\caption{공간 법칙 \(d\mathbf p/dt=\mathbf F\)과 그 시간 성분 짝인 일--에너지 관계.}\lectureframe{01:47:10}\end{figure}

이 방법은 최종적으로 식을 간결하게 만드는 것보다 중요합니다. 작용을 변분하기 전에 스칼라여야 한다고 요구했습니다. 그 결과 얻은 방정식은 나중에 대칭성을 억지로 보완하지 않고도 텐서 형태로 재구성할 수 있었습니다. 현대 장론에서는 이것이 일반적인 설계 원칙이 됩니다. 라그랑지안이 문제의 대칭성을 따르도록 구성하는 것입니다.

\section{국소성, 로런츠 불변성, 그리고 다음 대칭성}

상대론적 장론에서는 세 가지 원리가 거듭 등장합니다. 그중 두 가지는 이미 사용했고, 세 번째는 다음에 다루겠습니다.

\subsection{국소성}

국소성이란 입자나 장의 변화가 시공간에서 무한소로 가까운 곳에서 일어나는 일에 의해 좌우된다는 뜻입니다. 입자는 멀리 떨어진 무관한 값이 아니라 자신의 세계선 위에서 평가한 장에 반응합니다. 한 점의 장은 그 점의 근방에 있는 장과 그 미분에 반응하며, 임의의 먼 곳의 값이나 미래의 자료에 직접 반응하지 않습니다.

장에 대해서 작용은 다음과 같은 국소형을 취한다
\begin{equation}
S
=
\int d^3x\,dt\,
\mathcal L\!\left(\phi,\partial_\mu\phi\right).
\label{eq:sr6-local-action}
\end{equation}
적분은 시공간의 각 점에서 오는 기여를 하나씩 더한다. 일계 도함수는 이웃한 점들과의 비교를 나타내며, 변분 원리는 미분방정식을 도출한다. 서로 멀리 떨어진 점의 장들을 결합하는 비국소 곱을 생각할 수도 있지만, 그런 곱이 만드는 방정식에서는 이곳의 반응이 저곳의 조건에 직접 의존하게 된다. 여기서 부과하는 국소 구조는 그런 것이 아니다. 고계 도함수를 포함하는 작용도 가능하며 흥미로운 주제이지만, 지금의 구성을 위해서는 필요하지 않다.

\begin{figure}[htbp]
\centering
\includegraphics[width=0.90\textwidth]{lecture_06_figure_22.png}
\caption{장과 그 주변에서의 도함수로 이루어진 밀도의 시공간 적분으로 표현한 국소성.}
\lectureframe{01:51:00}
\end{figure}

\begin{classroomqa}
\audiencequestion{\(\partial_\mu\phi\)은 네 개의 도함수 성분을 나타내나요?}
\lecturerresponse{그렇습니다. 이는 시간 도함수와 \(\phi\)의 세 공간 도함수를 하나의 공변 4벡터로 모아 나타냅니다.\lecturetimestamp{01:52:25}}
\end{classroomqa}

\begin{classroomqa}
\audiencequestion{국소성 원리에서 \emph{가까운}이라는 말은 어떻게 정량화하나요?}
\lecturerresponse{도함수로 정량화합니다. 일계 도함수는 한 점에서의 장과 무한소만큼 이웃한 점들에서의 장 값을 비교합니다. 국소 라그랑지안은 같은 시공간 점에 있는 장들과 그 도함수 중 유한 개에 의존합니다.\lecturetimestamp{01:52:55}}
\end{classroomqa}

\begin{classroomqa}
\audiencequestion{방금 마친 입자 계산에 장의 국소성 개념은 어떻게 적용되나요?}
\lecturerresponse{입자 라그랑지안은 세계선 위에서 국소적입니다. 속도를 통해 이웃한 세계선 위치들에 의존하고, 입자의 현재 위치에서 평가한 전자기 퍼텐셜에도 의존합니다. 여기서 입자의 가속도는 멀리 떨어진 점의 장에 직접 의존하지 않습니다.\lecturetimestamp{01:53:31}}
\end{classroomqa}

\subsection{로런츠 불변성}

로런츠 불변성은 라그랑지안을 스칼라들로 구성하여 구현한다. 관찰자마다 이 스칼라들을 시간 성분과 공간 성분으로 다르게 분해할 수 있지만, 작용에는 동의한다. 따라서 변분으로 얻은 텐서 방정식은 모든 관성계에서 같은 형태를 갖는다.

\subsection{게이지 불변성}

게이지 불변성은 남은 원리다. 여기서는 이를 언급만 하고 아직 정의하지 않는다. 다음 단계에서 이 원리는 4퍼텐셜의 어떤 특징이 물리적으로 의미 있는지를 제약하고, 맥스웰 방정식을 만들어 내는 장 작용을 이끄는 역할을 한다.

\begin{figure}[htbp]
\centering
\includegraphics[width=0.90\textwidth]{lecture_06_figure_23.png}
\caption{세 가지 구성 원리: 국소성, 로런츠 불변성, 게이지 불변성.}
\lectureframe{01:55:40}
\end{figure}

이 원리들은 일반상대성이론, 양자전기역학, Yang--Mills 이론, 표준 모형에서 거듭 등장한다. 입자 종의 수와 같은 세부 사항은 이론마다 달라지지만, 국소성과 대칭성이 더 근본적인 틀을 제공한다. 따라서 실질적인 과제는 오일러--라그랑주 계산을 다시 따라가며 모든 지표와 묶인 항을 빠짐없이 설명할 수 있을 때까지 살펴보는 것이다. 그러면 게이지 불변성이 같은 전자기 작용에 대한 다음 제약 조건으로 등장한다.
